% ---------------------------------------------------------------------------
% Preamble: typography, math, theorem/exercise machinery, pedagogical boxes.
% ---------------------------------------------------------------------------
% Build with tectonic (XeTeX): native UTF-8, packages fetched on demand.
\usepackage{amsmath,mathtools,bm}         % math BEFORE newpxmath (symbol clashes)
\usepackage{amsthm}
\usepackage{newpxtext,newpxmath}          % Palatino-like text + matching math
\usepackage{microtype}
\usepackage[margin=1.15in]{geometry}
\usepackage{booktabs,tabularx,array,multirow}
\usepackage{enumitem}
\usepackage{graphicx}
\usepackage{xcolor}
\usepackage{tikz}
\usetikzlibrary{positioning,arrows.meta,fit,calc}
\usepackage[most]{tcolorbox}
\usepackage{listings}
\usepackage{caption}
\captionsetup{font=small,labelfont=bf}
\usepackage[colorlinks=true,linkcolor=linkblue,citecolor=linkblue,urlcolor=linkblue]{hyperref}
\usepackage[nameinlink,capitalise]{cleveref}

% ------------------------------------------------------------------- colors
\definecolor{linkblue}{HTML}{2B4C7E}
\definecolor{engbg}{HTML}{EEF3F8}
\definecolor{engframe}{HTML}{2B4C7E}
\definecolor{breakbg}{HTML}{FBEFEA}
\definecolor{breakframe}{HTML}{A84C32}
\definecolor{histbg}{HTML}{F4F1E8}
\definecolor{histframe}{HTML}{8A7B4D}
\definecolor{codebg}{HTML}{F7F7F7}
\definecolor{codeframe}{HTML}{999999}
\definecolor{kwcolor}{HTML}{2B4C7E}
\definecolor{strcolor}{HTML}{7A4A22}
\definecolor{comcolor}{HTML}{6B7B62}

% ------------------------------------------------------------------ listings
\lstdefinestyle{py}{
  language=Python,
  basicstyle=\ttfamily\footnotesize,
  keywordstyle=\color{kwcolor}\bfseries,
  stringstyle=\color{strcolor},
  commentstyle=\color{comcolor}\itshape,
  showstringspaces=false,
  columns=fullflexible,
  keepspaces=true,
  breaklines=true,
  frame=single,
  rulecolor=\color{codeframe},
  backgroundcolor=\color{codebg},
  framesep=4pt,
  xleftmargin=6pt, xrightmargin=6pt,
  aboveskip=8pt, belowskip=8pt,
}
\lstset{style=py}

% ------------------------------------------------------------------ theorems
\theoremstyle{plain}
\newtheorem{theorem}{Theorem}[chapter]
\newtheorem{proposition}[theorem]{Proposition}
\newtheorem{lemma}[theorem]{Lemma}
\newtheorem{corollary}[theorem]{Corollary}
\theoremstyle{definition}
\newtheorem{definition}[theorem]{Definition}
\newtheorem{example}[theorem]{Example}
\newtheorem{exercise}{Exercise}[chapter]
\theoremstyle{remark}
\newtheorem{remark}[theorem]{Remark}

% --------------------------------------------------------- pedagogical boxes
% Engineering consequence: where a piece of theory forces a code decision.
\newtcolorbox{engineering}[1][]{%
  enhanced, breakable, colback=engbg, colframe=engframe,
  boxrule=0.6pt, arc=1.5pt, left=6pt, right=6pt, top=4pt, bottom=4pt,
  title={Engineering consequence\ifblank{#1}{}{ --- #1}},
  fonttitle=\bfseries\small, coltitle=white}

% What breaks: the concrete failure that occurs if the lesson is skipped.
\newtcolorbox{whatbreaks}[1][]{%
  enhanced, breakable, colback=breakbg, colframe=breakframe,
  boxrule=0.6pt, arc=1.5pt, left=6pt, right=6pt, top=4pt, bottom=4pt,
  title={What breaks without this\ifblank{#1}{}{ --- #1}},
  fonttitle=\bfseries\small, coltitle=white}

% Origins: literature and history context.
\newtcolorbox{origins}[1][]{%
  enhanced, breakable, colback=histbg, colframe=histframe,
  boxrule=0.6pt, arc=1.5pt, left=6pt, right=6pt, top=4pt, bottom=4pt,
  title={Origins\ifblank{#1}{}{ --- #1}},
  fonttitle=\bfseries\small, coltitle=white}

% Code walkthrough: annotated excerpts from the real package.
\newtcolorbox{codewalk}[1][]{%
  enhanced, breakable, colback=white, colframe=codeframe,
  boxrule=0.6pt, arc=1.5pt, left=6pt, right=6pt, top=4pt, bottom=4pt,
  title={Code walkthrough\ifblank{#1}{}{ --- \texttt{#1}}},
  fonttitle=\bfseries\small, coltitle=white}

% ------------------------------------------------------------------- macros
\newcommand{\E}{\mathbb{E}}
\newcommand{\Prob}{\mathbb{P}}
\newcommand{\Var}{\operatorname{Var}}
\newcommand{\Cov}{\operatorname{Cov}}
\newcommand{\Corr}{\operatorname{Corr}}
\newcommand{\R}{\mathbb{R}}
\newcommand{\gauss}{\mathcal{N}}
\newcommand{\ind}[1]{\mathbf{1}\{#1\}}
\newcommand{\LSE}{\operatorname{LSE}}
\DeclareMathOperator*{\argmax}{arg\,max}
\DeclareMathOperator*{\argmin}{arg\,min}
\newcommand{\vx}{\bm{x}}
\newcommand{\vy}{\bm{y}}
\newcommand{\vz}{\bm{z}}
\newcommand{\vh}{\bm{h}}
\newcommand{\vpi}{\bm{\pi}}
\newcommand{\IC}{\mathrm{IC}}
\newcommand{\ICbar}{\overline{\mathrm{IC}}}
\newcommand{\ICIR}{\mathrm{ICIR}}
\newcommand{\NLL}{\mathrm{NLL}}
\newcommand{\code}[1]{\texttt{#1}}

% Exercise list at chapter end
\newcommand{\exercisehead}{\section*{Exercises}\addcontentsline{toc}{section}{Exercises}}
